\documentclass{book}
\usepackage[utf8]{inputenc}
\usepackage[T1]{fontenc}
\usepackage{geometry}
\usepackage{graphicx}
\usepackage{float}
\usepackage{hyperref}
\usepackage[english,indonesian]{babel}

\title{Optimasi Pemrosesan Template LaTeX Menggunakan Pendekatan Berbasis AI}
\author{Academic Support Agent}
\date{\today}

\begin{document}
\maketitle

\begin{abstract}

\end{abstract}

\chapter{PENDAHULUAN}


\chapter{TINJAUAN PUSTAKA}


% Original content (commented out):
% \documentclass[12pt,a4paper,openany]{book} %el comando "openany" sirve que cuando se hace un Chapter, dejará una hoja en blanco.
\usepackage[utf8]{inputenc}
\usepackage[T1]{fontenc}
\usepackage[spanish,es-tabla]{babel}
\usepackage{amsmath}
\usepackage{amsfonts}
\usepackage{amssymb}
\usepackage{graphicx}
\usepackage{wasysym}
\usepackage{lscape}
\usepackage[colorlinks=true, linkcolor=black,
citecolor=black, urlcolor=black,hidelinks]{hyperref}
\usepackage{float}
\usepackage{booktabs}
\usepackage{setspace} %por si el titulo es mas largo y el interlineado
\usepackage{xcolor} %para darle color a las lineas del titulo.
\usepackage{lipsum} %crea texto ficticio del texto Lorem Ipsum
%\definecolor{azul}{rgb}{0.17, 0.4, 0.69} %definir color de lineas.
\usepackage{natbib}
\usepackage[left=2.54cm,top=2.54cm,right=2.54cm,,bottom=2.54cm]{geometry} %para cambiar la configuracion de las medidas a APA se usa "Control + R".
\newcommand{\grad}{$^{\circ}$}
%\renewcommand{\baselinestretch}{1.5} %interlineado 1.5
\author{Jorge Vicente Niño Mocarro} %integrantes
\title{Tesis} %titulo del trabajo

\begin{document}

	%FRONTMATTER
	%---------------------------------------------------
	\frontmatter %numeros romanos
	\include{caratula}
	\include{jurado}
	%DEDICATORIA
	\include{agradecimientos}
	%ÍNDICE O TABLA DE CONTENIDO-------------------
	\tableofcontents
	%\cleardoublepage
	\addcontentsline{toc}{chapter}{Lista de figuras} % para 						que aparezca en el indice de contenidos
	\listoffigures % indice de figuras
	%\cleardoublepage
	\addcontentsline{toc}{chapter}{Lista de tablas} % para que 						aparezca en el indice de contenidos
	\listoftables % indice de tablas
	
	 % si no queremos que añada la palabra "Capitulo"
	%\addcontentsline{toc}{chapter}{Agradecimientos} % si queremos que aparezca en el índice
	%\markboth{AGRADECIMIENTOS}{AGRADECIMIENTOS} % encabezado
	
	
	
	%MAINMATTER
	%--------------------------------------------------	
	\mainmatter %numeros arabicos.
	%CAPITULO I:PLANTEAMIENTO DEL PROBLEMA
	\include{capitulos/planteamiento}
	
	%CAPITULO II:MARCO TEORICO
	\include{capitulos/marcoteorico}
	
	%CAPITULO III: DISEÑO METODOLOGICO
	\include{capitulos/diseño}
	
	%CAPITULO IV:ACTIVIDADES Y RECURSOS
	\include{capitulos/actividades}
	%--------------------------------------------------
	
	
	
	%BACKMATTER
	%---------------------------------------------------
	\backmatter
	%BIBLIOGRAFÍA.
	%\cleardoublepage
	\addcontentsline{toc}{chapter}{Bibliografía}
	\bibliographystyle{apalike}
		\bibliography{capitulos/biblio}
		
	%---------------------------------------------------	
	
\end{document}


%NO OLVIDAR CITAR!

\end{document}